\section{ReproBench}
\label{sec:reprobench}

ReproBench evaluates whether an autonomous agent can reproduce the central
empirical result of a machine learning paper from whatever artifacts its
authors actually released. Each of the 100 papers in the benchmark poses one
task. The agent must recover a single pinned result within a stated tolerance.
This section defines the task (Section~\ref{subsec:task}) and describes how
the benchmark was constructed (Section~\ref{subsec:construction}).

\subsection{Task definition}
\label{subsec:task}

Each paper carries one reproduction target, which we call the \emph{match
target} and fix before any agent runs. The match target is the cheapest
configuration the authors report that still establishes the paper's central
empirical claim, together with the value that configuration produces. Papers
usually headline their largest model or their full benchmark suite, but the
same claim typically rests on a smaller supporting experiment, and we pin that
one. Pinning the smallest supporting configuration bounds the compute a task
can demand while still exercising the paper's method end to end.

% TODO(data): one eval-100 row currently has a null match_bar_kind; the
% "remaining 3" count below assumes it resolves to magnitude or range.
A match target records the configuration that produces the value, the metric,
the reported value, the scope the value is measured over, and the shape of the
bar a run must clear. For 67 of the 100 papers the bar is a point estimate.
For 18 the run must beat a stated baseline, for 12 it must clear a reported
threshold, and for the remaining 3 it must recover a magnitude or fall inside
a reported range. A point estimate matches when the run lands within the
authors' stated tolerance, or within 5\% relative of the reported value when
the paper states none.

As a running example, SharpZO (arXiv 2506.20990) claims that a two-stage
zeroth-order method improves prompt tuning for vision-language models using
forward passes alone. Its match target pins a test accuracy of 79.42\% on the
EuroSAT test set under a 16-shot CLIP ViT-B/16 configuration the paper
reports. A run matches when it produces that accuracy from an experiment it
actually executed, within the default 5\% relative tolerance.

The agent receives the paper, the metric, the target value, the scope, and
the bar shape. It does not receive the configuration. The agent must
therefore design the experiment itself, deciding from the paper and the
released artifacts what to run, retrain, or reimplement. Because the agent
knows the number it must hit, a bare value match proves nothing. Grading is
provenance-based. An LLM auditor traces every claimed result to execution
evidence in the run directory, and deterministic checks force the score to
zero when a number was echoed from the paper rather than measured.

\subsection{Benchmark construction}
\label{subsec:construction}

We construct ReproBench from the 3,414 NeurIPS 2025 submissions that carry an
arXiv identifier, collected in the \emph{ai-conferences}
corpus~\citep{neurips2025corpus}. Construction proceeded in three stages.

\textbf{Stage I: Artifact-verified classification.} For each submission we
retrieved the raw \LaTeX{} source together with any OpenReview supplementary
material, since some authors ship code in the supplement rather than on
GitHub. A classifier agent built on MiniMax-M2.7~\citep{minimax2026m2} read
each paper, extracted its central empirical claim, pinned its match target,
and audited four artifact signals: released code, released datasets, released
model weights, and whether the required datasets are standard public
benchmarks. The audit trusts tool calls over prose. An artifact counts as
available only when a tool call opens it. An artifact the paper promises but
never publishes counts as missing, as does one behind a paywall, a sign-up
form, or a permission request. We ran the classifier on a 1,000-paper sample
of the pool, and it fully verified 687 papers. Of these, 34 are non-empirical
and 33 have their required artifacts blocked, leaving 620 papers eligible for
evaluation.

Artifact availability sets each paper's tier (Table~\ref{tab:tiers}). The
tier records what the agent must rebuild, not the difficulty of the
underlying science. Compute cost is tracked separately, so a tier records
access, not scale. The tiers also describe the venue rather than a curated
corner of it. Among the 687 verified papers, 33\% fall in Easy, 44\% in
Medium, and 13\% in Hard, while 5\% are blocked outright and 5\% are
non-empirical.

\begin{table}[t]
\caption{The three artifact-availability tiers. Each tier states what the
agent must supply for itself and therefore how much of the result it must
rebuild.}
\label{tab:tiers}
\begin{center}
\begin{tabular}{@{}lll@{}}
\toprule
Tier & Artifacts available & Agent's job \\
\midrule
Easy & Code, data, and weights & Run or evaluate the released artifacts. \\
Medium & Code and data, no weights & Retrain the match-target model, then evaluate. \\
Hard & Data only, no code & Reimplement the method from the paper. \\
\bottomrule
\end{tabular}
\end{center}
\end{table}

\textbf{Stage II: Compute capping and stratified selection.} We estimated the
compute each match target requires in H100-hours, taking the paper's reported
hardware and wall-clock time where stated and deriving the estimate from the
pinned configuration otherwise. We dropped every paper whose estimate exceeds
192 H100-hours, which caps what a single reproduction can demand. From the
remainder we sampled a 200-paper audit pool of 67 Easy, 67 Medium, and 66
Hard papers, stratified across four compute bands (0--8, 8--32, 32--96, and
96--192 H100-hours).

\textbf{Stage III: Human audit.} We audited every paper in the pool by hand,
opening the evidence behind each availability call and re-deriving each
compute estimate. Papers failed the audit for two reasons. Six papers had
required artifacts that turned out to be inaccessible, and fifteen depended
on paid third-party APIs, which we exclude so that reproduction cost stays
bounded and GPU-denominated. Rejected papers were replaced from the audited
pool. Only 17 of the 100 shipped compute estimates could be taken directly
from the paper. We derived the rest from the pinned configuration, echoing
the finding that compute is systematically
under-reported~\citep{gundersen2018state}.

Through these stages the 3,414 submissions narrow to the released benchmark.
It holds a 100-paper evaluation split of 34 Easy, 33 Medium, and 33 Hard
papers, and a disjoint 14-paper development split. After adjudication the
evaluation split totals 1,651 H100-hours, its compute bands hold 58, 25, 14,
and 3 papers, and exactly one paper saturates the 192-hour cap. Of the 100 papers, 94
first appeared as preprints in 2025, which limits the exposure of their
artifacts in current training corpora.
